%% Die Sprachdatei von prag für de-DE (ADR-0025).

%% ---------------------------------------------------------------- Zahlen
%% Ab wie vielen Stellen gruppiert wird, wählt die Vorlage (ADR-0021).
\sisetup{
  output-decimal-marker = {,},
  group-separator       = {.},
}
\newcommand{\zeichenEUR}{€}
\newcommand{\zeichenUSD}{\$}
\newcommand{\zeichenGBP}{£}
\newcommand{\zeichenPLN}{PLN}
\newcommand{\zeichenCZK}{CZK}
\newcommand{\zeichenRON}{RON}
\newcommand{\zeichenSEK}{SEK}
\newcommand{\zeichenNOK}{NOK}
\newcommand{\zeichenCHF}{CHF}
\newcommand{\waehrungszeichen}{\csname zeichen\waehrung\endcsname}
\newcommand{\betrag}[1]{\num{#1}~\waehrungszeichen}
\newcommand{\eurobetrag}[1]{\begingroup\def\waehrung{EUR}\betrag{#1}\endgroup}
\newcommand{\prozent}[1]{\num{#1}~\%}

%% ------------------------------------------------------------- Seiten
\newcommand{\wortseitevon}[2]{Seite #1 von #2}
\newcommand{\wortfortsetzung}{Fortsetzung auf der nächsten Seite}

%% ------------------------------------------------------------ Kopfangaben
\newcommand{\wortrechnungsnr}{Rechnungsnr.:}
\newcommand{\wortkundennr}{Kundennr.:}
\newcommand{\wortauftragsnr}{Unsere Auftragsnr.:}
\newcommand{\wortdatum}{Datum:}
\newcommand{\wortvereinnahmtam}{Vereinnahmt am:}
\newcommand{\wortreferenz}{Ihre Referenz:}

%% ----------------------------------------------------- Positionen, Summen
\newcommand{\wortpos}{Pos.}
\newcommand{\wortbezeichnung}{Bezeichnung}
\newcommand{\wortmenge}{Menge}
\newcommand{\worteinheit}{Einheit}
\newcommand{\worteinzelpreis}{Einzelpreis}
\newcommand{\wortrabatt}{Rabatt}
\newcommand{\wortgesamt}{Gesamt}
\newcommand{\wortzwischensumme}{Zwischensumme (netto)}
%% Satz und, bei mehreren Steuergruppen, ihre Grundlage – sonst ist sie leer.
\newcommand{\wortsteuer}[2]{Umsatzsteuer #1\ifblank{#2}{}{ auf #2}}
\newcommand{\wortgesamtbetrag}{Gesamtbetrag}

%% --------------------------------------------------------------------- Fuß
%% Ohne Doppelpunkt, den setzt die Vorlage hinter jede Beschriftung.
\newcommand{\wortustid}{USt-IdNr.}
\newcommand{\worttelefon}{Tel.}
\newcommand{\wortsteuernummer}{Steuernummer}
\newcommand{\wortregistergericht}{Registergericht}
\newcommand{\wortgeschaeftsfuehrung}{Geschäftsführung}
